\scanpage{38}
% archive sheet number: 37
\deleted{\unsure{Et} \inserted{Alors} $F_M$ est donné par les entiers $m_i$ \inserted{et par $\varphi_{f-1}$} \unsure{mais qui} vérifient donc \unsure{$=$}, les $\unsure{m_i,m_{i-2}|m_i,m_{i-1}|m_{i+1}-m_i}$}

On a choisi les isom de telle façon que chaque $\varphi_{\unsure{\nu|v}}$ ($0\leq\nu\leq f-2$) soit de la forme
\[
\varphi_\nu=\pi^{r_{\deleted{\unsure{0}}\nu}}\mathrm{Id}_\Omega
\qquad (r_\nu\geq 0\deleted{\unsure{)}},\;0\leq\nu\leq f-2)
\]
et $\varphi_{f-1}$ est donné alors par
\[
\varphi_{f-1}=\lambda\,\mathrm{id}_\Omega \qquad \lambda\in A,\;\inserted{\text{s'écrit}}\ \lambda=\pi^{r_{f-1}}u,\;u\in A^*
\]

\deleted{\unsure{La co}} On a alors
\[
F_M^f=\pi^{r_0+\cdots+r_{f-2}+r_{f-1}}u=\varpi\qquad u\in A^*
\]

\deleted{On aura}
\[
\deleted{M/F_M M\simeq A/\pi^{r_{f-1}}A+A/\pi^{r_{\unsure{0|i}}}A+\ldots+A/\pi^{r_{f-2}}A}
\]

La condition $(*)'$ s'écrit

\deleted{\unsure{donc} \unsure{donc|$d_K$} $\pi$}

\[
\begin{aligned}
&p\ \text{divisible par }\pi^{r_\nu}\ \text{pour tt }0\leq\nu\leq f-2,\\
&\hphantom{p\ \text{divisible }}\text{par }u\pi^{r_{f-1}}
\end{aligned}
\]
i.e.
\[
\left\{\begin{aligned}
&\boxed{\deleted{\unsure{0\leq}}r_\nu\leq e}\quad\text{pour }0\leq\nu\leq f-1\\
&\qquad\boxed{\deleted{\unsure{r_{f-1}+v(u)\leq e}}}
\end{aligned}\right.
\]
On a
\[
M/F_M M=A/u\pi^{r_{f-1}} A\unsure{+|-}A/\pi^{r_0}A+\ldots+A/\pi^{r_{f-2}}A
\]
est de \deleted{\unsure{dimension}} longueur
\[
\begin{aligned}
&\boxed{\dim G_0=\deleted{\unsure{v(u)+}} \inserted{\unsure{\operatorname{long} A(\varpi A)\mathrel{\overset{?}{=}}|\operatorname{long} A/(\varpi A)\mathrel{\overset{?}{=}}}}\ r_0+\ldots+r_{f-2}+\unsure{(}r_{f-1}\deleted{\unsure{+v(u)}})}\\
&\hspace{7em}(\leq fe=h)
\end{aligned}
\]

\deleted{Supposons maintenant}

Donc on trouve
\[
\boxed{\begin{aligned}
\dim G_0=1&\Longleftrightarrow F_M^f\ \deleted{\unsure{\text{une}}}\inserted{=\varpi_M,\;\varpi\ \text{une}}\ \text{\uline{uniformisante} de }A\\[-2pt]
&\qquad\text{(qu'on peut supposer alors }=\pi,\ \text{quitte à changer }\pi\,!\,)\\
&\Longleftrightarrow\text{tous les }r_\nu\text{ sauf un seul }\inserted{r_{\nu_0}}\text{ sont nuls},\\
&\qquad r_{\nu_0}=1,\quad u\text{ une unité}\\
&\hspace{9em}(u=1\text{ si on veut }\varpi=\pi)
\end{aligned}}
\]

\unsure{Ce|G}
